\section{Uniform rational approximation of polynomial inverses}
\label{app:inverse}

This appendix proves the approximation interface used in
\cref{sec:accuracy}. Polynomial inverse approximation and algebraic-series
computation have established algorithms
\cite{Farouki2000,Walsh2000}. The purpose here is to give the complete
uniform error and rational encoding bounds, including flat points of a
strictly increasing marginal. The analytic ingredients are the classical
holomorphic inverse, Cauchy estimates, and Rouch\'e theorem
\cite{DLMFComplex}.

\begin{lemma}[Effective increment and inverse bounds]\label{lem:inverse-modulus}
Let $g\in\Q[t]$ have degree $D\ge1$, be strictly increasing on $[0,1]$,
and satisfy $g(0)=0$, $g(1)=1$. For $[a,a+h]\subseteq[0,1]$,
\begin{equation}\label{eq:inverse-increment}
             g(a+h)-g(a)\ge\frac12\left(\frac{h}{2D}\right)^D.
\end{equation}
Consequently the clipped inverse changes by less than $e$ whenever its
argument changes by less than $\frac12(e/(2D))^D$, for $0<e\le1$.
For an unnormalized marginal with endpoint difference $G>0$, multiply
this argument tolerance by $G$.
\end{lemma}
\begin{proof}
For $h>0$ put $v=g(a+h)-g(a)$. Interpolate $g-g(a)$ at
$a+jh/D$, $j=0,\ldots,D$. All node values lie in $[0,v]$.
On $[0,1]$ every factor in the numerator of a Lagrange basis polynomial
has absolute value at most one, and its denominator has magnitude
$(h/D)^D j!(D-j)!$. Thus
\[
 |g(t)-g(a)|\le v(D/h)^D\sum_{j=0}^D\frac1{j!(D-j)!}
              \le v(2D/h)^D.
\]
Applying this at $t=0,1$ gives $1\le2v(2D/h)^D$. The case $h=0$
is immediate. An inverse displacement of at least $e$ contains a
response interval of length $e$, so requires at least the displayed
argument displacement. If clipping occurs, the argument interval only
gets longer. Positive affine normalization proves the last assertion.
\end{proof}

\begin{theorem}[Rational inverse approximation]\label{thm:inverse-approximation}
For a densely encoded rational polynomial $g$ strictly increasing on
$[0,1]$, and a positive rational tolerance $\eta$, one can construct
rational breakpoints and rational polynomial branches approximating its
clipped inverse uniformly to error $\eta$. Each branch is valid on its
closed interval, with constant branches outside $[g(0),g(1)]$.
Construction time, number of branches, degrees, and rational encoding
lengths are polynomial in input bits, numerical degree, and
$\max\{0,\log(1/\eta)\}$. The encoding of $\eta$ is part of the input;
equivalently one can request $\eta=2^{-B}$ at cost polynomial in $B$.
No positive lower bound on $g'$ is required.
\end{theorem}
\begin{proof}
Normalize to $g(0)=0,g(1)=1$ by a rational affine change of target.
Replace $\eta$ by $\min\{\eta,1/2\}$, and put
\begin{equation}\label{eq:inverse-bad-tolerance}
 e=\eta/4,\qquad \delta_0=\tfrac14(e/(2D))^D,\qquad
 \rho=\text{largest dyadic number at most }
          \min\{\delta_0/[4(D+1)],1/16\}.
\end{equation}
Their bit lengths are polynomial in the stated parameters. An interval
of target length at most $\delta_0$ has inverse oscillation less than
$e$, by \cref{lem:inverse-modulus}.

\emph{Critical values, including nonreal ones.}
The real parts of all finite complex critical values are the real $s$
satisfying
\begin{equation}\label{eq:critical-projection}
 \exists u,v:\quad
 \Re g'(u+\mathrm i v)=\Im g'(u+\mathrm i v)=0,\qquad
                         s=\Re g(u+\mathrm i v).
\end{equation}
The formula has three real variables and polynomials of degree at most
$D$ with polynomial-size dense expansions. Fixed-dimensional elimination
and univariate isolation from \cref{subsec:tools} compute this finite
set, of cardinality at most $D-1$, in polynomial bit time. For $D=1$ it
is empty. Repeated critical points or coincident real parts cause no
problem. Keeping only real critical values would be insufficient for the
analytic argument below.

Retain the real parts in $[-1,2]$, bracket each in a rational interval
of width at most $\rho$, and pad that interval by $\rho$ at both ends.
Add $[-\rho,\rho]$ and $[1-\rho,1+\rho]$, intersect the union with
$[0,1]$, and merge overlapping intervals. Denote the result by $B_0$.
The sum of untrimmed lengths is at most
$3(D-1)\rho+4\rho\le\delta_0$. Each merged component therefore has
length at most $\delta_0$. On it use a constant rational approximation,
within $e$, to the inverse at its rational midpoint. Monotone bisection
in response space computes that constant in polynomial time. Its error
on the entire component is less than $2e$. This also treats endpoints
and interior derivative zeros.

For a complementary gap $[a,b]$, define
\[
                    d(t)=\min\{t-a+\rho,b-t+\rho\}.
\]
Every critical real part is at distance at least $d(t)$ from $t$.
For retained parts, padding places them at least $\rho$ behind the
exposed endpoint of their merged component. Omitted parts lie outside
$[-1,2]$, at distance greater than one, whereas
$d(t)\le1/2+\rho<1$.

\emph{Rational analytic panels.}
Split each gap at $m=(a+b)/2$. Starting at $x=a$, set
\[
                  y=\min\{m,x+(x-a+\rho)/32\}
\]
and take $[x,y]$ as the next panel; reflect this construction on the
right half. Before the final panel, $x-a+\rho$ grows by $33/32$.
Thus each gap has $O(1+\log(1/\rho))$ panels, all with polynomial
rational bit length. There are at most $D+1$ gaps. At a panel midpoint
$\tau$, put $d=d(\tau)$; the panel half-width is at most $d/64$.

Writing $g(z)=\sum a_jz^j$, the rational number
$L_g=\max\{1,\sum j|a_j|\}$ bounds $g'$ in absolute value on
$[0,1]$. Bisection of the inverse bracket to width $d/(64L_g)$
provides a rational midpoint $z_0\in(0,1)$ such that
$|g(z_0)-\tau|\le d/64$. Set $t_0=g(z_0)$ and $R=d/4$.
Every complex critical value has distance at least $63d/64$ from
$t_0$. Hence a neighborhood of the closed radius-$R$ disk contains
no critical value. Every panel target satisfies
\begin{equation}\label{eq:inverse-panel-ratio}
                       |t-t_0|\le d/32\le R/8.
\end{equation}

For completeness, the inverse branch exists on that entire disk, not
merely on its real diameter. A polynomial is proper: the modulus of its
leading term forces $|g(z)|\to\infty$ as $|z|\to\infty$.
Away from critical values each preimage is simple, and the local
holomorphic inverse theorem gives disjoint inverse neighborhoods for
its finitely many roots. Properness ensures that these neighborhoods
exhaust preimages after shrinking the target neighborhood; otherwise a
sequence of additional preimages has a bounded convergent subsequence,
contradicting local uniqueness. Thus the restriction is a finite
covering. Path lifting over the simply connected disk gives one inverse
on each connected component, with no monodromy. It is holomorphic by
local inversion. Select the component through $z_0$. Along the real
segment to every panel target this branch equals the real increasing
inverse by uniqueness and continuity. In particular $g'(z_0)\ne0$.

\emph{Taylor approximation and rational arithmetic.}
On the disk $|t|<2$. The elementary polynomial root bound gives
\[
 |Z(t)|\le M_g:=2+\frac{2+\sum_{j<D}|a_j|}{|a_D|}.
\]
Indeed, a root of $a_Dz^D+\sum_{j<D}\widetilde a_jz^j$ has
modulus at most $1+\max_{j<D}|\widetilde a_j/a_D|$, and the
constant coefficient changes by less than two. Cauchy's estimate for
$Z(t_0+v)=z_0+\sum_{n\ge1}c_nv^n$ yields
$|c_n|\le M_gR^{-n}$. With
$q\ge\lceil\log_2(4M_g/\eta)\rceil$, the tail on the panel is
at most $M_g2^{-q}\le\eta/4$, using the weaker ratio $1/2$ in
\eqref{eq:inverse-panel-ratio}.

All coefficients are rational, since the expansion center is a
\emph{rational response} rather than an approximate irrational inverse.
Write
\[
 g(z_0+u)-t_0=Q^{-1}\sum_{h=1}^D A_hu^h,
 \qquad Q\in\Z_{>0},\ A_h\in\Z,\ A_1\ne0.
\]
Formal substitution gives
\begin{equation}\label{eq:inverse-reversion}
 \begin{split}
 c_1&=Q/A_1,\\
 c_n&=-\frac1{A_1}\sum_{h=2}^{\min(D,n)}A_h[v^n]
                       \left(\sum_{j=1}^{n-1}c_jv^j\right)^h.
 \end{split}
\end{equation}
A factor containing $c_n$ cannot contribute to degree $n$ in a power
$h\ge2$, which proves the triangular recurrence. Inductively
\begin{equation}\label{eq:inverse-denominator}
                    c_n=Q^nN_n/A_1^{2n-1},\qquad N_n\in\Z.
\end{equation}
A product of $h$ earlier coefficients with indices summing to $n$
has denominator $A_1^{2n-h}$ after factoring $Q^n$. The additional
division by $A_1$ and conversion to the displayed denominator multiplies
its integer numerator by $A_1^{h-2}$. Signs are irrelevant.

Translation of $g$ to $z_0$ has polynomial encoding; so do $Q,A_h$.
The denominator in \eqref{eq:inverse-denominator} and Cauchy's magnitude
bound give polynomial numerator and denominator bit lengths, because
$q$, $\log M_g$, and $\log(1/R)$ are polynomially bounded. The same
holds for intermediate truncated convolutions: each term has at most
$q$ earlier factors, and at most $2^{q-1}$ compositions contribute to
any fixed degree before the polynomially many sums over $h$.
Using these common denominators keeps all arithmetic polynomial in bit
cost. Truncated polynomial multiplication uses polynomially many
operations, without enumerating the compositions. Ordinary-power
expansion of the Taylor polynomials preserves the bounds.

The constant pieces on $B_0$, Taylor pieces on the gaps, and exact
clipped constants cover the line. Adjacent branches each satisfy the
error bound at their shared endpoint; they need not agree exactly.
A half-open convention defines a single approximating function if needed.
Undoing the target normalization proves the theorem.
\end{proof}

\begin{proposition}[Sharper positive-coefficient construction]
\label{prop:positive-inverse}
If $g(z)=\sum_{j=1}^P a_jz^j$ with $a_j\ge0$ and $g(1)>0$,
the approximation in \cref{thm:inverse-approximation} can use
$O(P^2\log(1/\eta))$ nonconstant branches, each of degree
$O(\log(1/\eta))$, independently of coefficient conditioning in these
two counts. Coefficient bit lengths still depend polynomially on the
input encoding. Numerical, rather than sparsely encoded binary, $P$
is the complexity parameter.
\end{proposition}
\begin{proof}
Normalize $g(1)=1$. Termwise,
\[
 z^P\le g(z)\le z,\quad 0\le g'(z)\le P,\quad
                g(z)\le zg'(z)\le Pg(z)\quad(0<z\le1).
\]
Take $m=\lceil\log_2(1/\eta)\rceil\ge1$ after reducing
$\eta$ to at most $1/2$, and $K_0=Pm$. Below $2^{-K_0}$ the
constant zero has error at most $2^{-m}$; above one use the exact
constant one.

For real $0<z_0\le1$, put $t_0=g(z_0)$,
$r_z=z_0/(4P)$ and $R_t=t_0/(8P)$. Nonnegative coefficients give,
for complex $|u|\le r_z$,
\[
 |g'(z_0+u)-g'(z_0)|
 \le g'(z_0)((1+1/(4P))^{P-1}-1)<g'(z_0)/3.
\]
Here $e^{1/4}<4/3$ follows by termwise comparison with a geometric
series. Integrating along the segment gives
\[
 |g(z_0+u)-t_0-g'(z_0)u|\le g'(z_0)|u|/3.
\]
On $|u|=r_z$, the shift $|t-t_0|\le R_t$ is at most
$g'(z_0)r_z/2$. The nonlinear remainder plus shift is therefore
at most $5/6$ of the linear comparison magnitude. Rouch\'e gives
exactly one root in the response disk; the derivative bound makes it
simple. Local inverses agree by uniqueness, and the strict boundary
margin extends them to a neighborhood of the closed target disk.
The resulting branch satisfies $|Z(t)|<2$ and agrees with the
positive real inverse on real targets in $[0,1]$.

Split each dyadic interval $[A,2A]$, $A=2^{-(k+1)}$,
$k=0,\ldots,K_0-1$, into $32P$ equal rational panels. At midpoint
$\tau$ the panel half-width is at most $\tau/(64P)$.
Bisection to response width $\tau/(64P^2)$ produces rational
$z_0$ with $|t_0-\tau|\le\tau/(64P)$ and
$t_0\ge63\tau/64$. Thus on its panel
\[
             |t-t_0|/R_t\le16/63<1/2.
\]
The bisection takes $K_0+O(\log P)$ steps. Use the exact recurrence
\eqref{eq:inverse-reversion}. Cauchy's estimate is now
$|c_n|\le2R_t^{-n}$, so truncating at $q=m+3$ gives tail
$2^{1-q}\le\eta/4$. The denominator and intermediate-bit proof
above applies unchanged. There are $32P^2m$ panels and the claimed
degree. This argument uses coefficient nonnegativity essentially in
the derivative-disk estimate; it is not used for signed marginals.
\end{proof}
